\documentclass[12pt]{article}
\usepackage[utf8]{inputenc}
\usepackage[T1]{fontenc}
\usepackage{hyperref}
\title{Lower Ta'oih Language Resources for K-12 Curriculum Development}
\author{Compiled by Claude Sonnet 5 (claude-sonnet-5) for Andr\'as Kornai}
\date{}
\begin{document}
\maketitle

\section{Introduction}
Lower Ta'oih (ISO 639-3: tto; Glottocode lowe1395) is an Austroasiatic language of Laos. It belongs to the Ta'oih (Ta'Oi) group, along with Upper Ta'oih, Ir and Ong. Sources disagree on whether these varieties are separate languages or a dialect continuum. Searches turned up no resource specific to Lower Ta'oih.

\section{Primary Educational Resources}
The searches found no primers, textbooks or school programmes for Lower Ta'oih.

\section{Linguistic and Reference Materials}
\begin{itemize}
\item Wikipedia article on the Ta'Oi language (the wider group): \url{https://en.wikipedia.org/wiki/Ta%27Oi_language}
\item Omniglot page on Ta'Oi. It describes a Lao-script orthography devised by Feikje van der Haak of SIL for the Kataang dialect: \url{https://omniglot.com/writing/taoi.htm}
\item van der Haak, F. (1993), ``Ta'uaih phonology and orthography: a preliminary statement'', \emph{Mon-Khmer Studies} 22: 257--264, as catalogued by Glottolog under Upper Ta'oih. The catalogue entry is: \url{https://glottolog.org/resource/reference/id/154395}
\item Ethnologue subgroup listing: \url{https://ethnologue.com/subgroup/1759/}
\end{itemize}

\section{Religious and Community Resources}
The only item found is a Joshua Project profile for Lower Ta Oi in Laos. It is a demographic profile, not a language resource.
\begin{itemize}
\item \url{https://joshuaproject.net/people_groups/19360}
\end{itemize}

\section{Digital and Audio Resources}
No audio, video or corpus resources were found.

\section{Curriculum Development Context}
The language is stable and used by the whole community, but I found no written materials specific to it. The only orthography found is a Lao-script scheme for the Kataang dialect. Whether it applies to Lower Ta'oih is not established here. Any curriculum would need new materials developed with the community.
\end{document}